% Generated from locale/tr/content/lambda-calculus/syntax/abbreviated-syntax.tex; do not hand-edit.


\olfileid[tr]{lam}{syn}{abb}
\olsection{Kısaltılmış Sözdizim}

\olref[trm]{defn:term} içinde tanımlanan terimleri yazmak bazen hantaldır;
bu nedenle daha özlü bir sözdizim kullanmak yararlıdır. Elbette özlü
gösterimdeki terimlerin de tek türlü okunabildiğinden emin olmalıyız. Yaygın
olarak kullanılan \emph{kısaltılmış terimler} gösterimi şöyledir.

\begin{enumerate}
\item Parantezler atıldığında uygulama soldan sağa yapılır. Örneğin $M$,
  $N$, $P$ ve~$Q$ terimse $MNPQ$, $(((MN)P)Q)$ terimini kısaltır.
\item Yine parantezler atıldığında lambda soyutlamasına olabilecek en geniş
  kapsam verilir. Örneğin $\lambd[x][MNP]$, $(\lambd[x][MNP])$ olarak
  okunur.
\item Tek bir lambda birden çok değişkeni soyutlamak için kullanılabilir.
  Örneğin $\lambd[xyz][M]$,
  $\lambd[x][\lambd[y][\lambd[z][M]]]$ ifadesinin kısaltmasıdır.
\end{enumerate}

Örneğin
\[
\lambd[xy][xxyx \lambd[z][xz]]
\]
ifadesi
\[
(\lambd[x][(\lambd[y][((((xx)y)x)(\lambd[z][(xz)]))])]).
\]
ifadesini kısaltır.

\begin{prob}
$\lambd[g][(\lambd[x][g (x x)])
  \lambd[x][g (x x)]]$ kısaltılmış terimini açın.
\end{prob}

